% Generated draft: requires mathematical and source review.
\documentclass[11pt]{article}
\usepackage[margin=25mm]{geometry}
\usepackage{fontspec}
\setmainfont{Latin Modern Roman}
\usepackage{amsmath,amssymb,mathtools,unicode-math}
\setmathfont{Latin Modern Math}
\usepackage{enumitem,needspace,xurl,hyperref}
\setlength{\parindent}{0pt}
\setlength{\parskip}{5pt}
\setlength{\emergencystretch}{3em}
\newcommand{\N}{\mathbb{N}}
\newcommand{\Z}{\mathbb{Z}}
\newcommand{\Q}{\mathbb{Q}}
\newcommand{\R}{\mathbb{R}}
\newcommand{\C}{\mathbb{C}}
\newcommand{\F}{\mathbb{F}}
\newcommand{\A}{\mathbb{A}}
\newcommand{\PP}{\mathbb{P}}
\newcommand{\E}{\mathbb{E}}
\newcommand{\eps}{\varepsilon}
\newcommand{\dd}{\,\mathrm{d}}
\begin{document}
\section*{UnsolvedMath ID: 6200090}
\textbf{Source problem code:} catalogue entry
\subsection{Definitions and mathematical statement}
A graph of groups assigns groups to vertices and edges, with injective maps from each edge group to its incident vertex groups. A complex of groups generalizes this to higher-dimensional cells: local groups are assigned to cells, with injective face maps and, in the general case, twisting elements satisfying compatibility and cocycle identities. A developable complex of groups is one realized by a group acting on a simply connected complex with those local stabilizers. Developability is an additional issue in higher dimensions; it must not be assumed from injectivity of the local maps alone.

The requested generalization is a Bestvina--Feighn-style combination theorem for the fundamental group \(\pi_1(\mathcal G)\) of a complex of groups \(\mathcal G\). In the hyperbolicity interpretation, the objective is to formulate geometric and algebraic hypotheses on the local groups, face embeddings, higher-dimensional gluing, and development that imply
\[
\pi_1(\mathcal G)\text{ is a Gromov-hyperbolic group},
\]
meaning that a Cayley graph for a finite generating set has uniformly thin geodesic triangles. The theorem should specialize appropriately to the graph-of-groups combination result when the underlying complex is one-dimensional. It is not an assertion that arbitrary complexes of hyperbolic groups have hyperbolic fundamental groups.

The source does not choose a version of the Bestvina--Feighn theorem, a class of base complexes, developability assumptions, metric realization, or an analogue of the quasi-isometric embedding and flaring hypotheses used in combination settings. Determining suitable assumptions is part of the programme, so no universal sufficient list or optimal criterion is invented here. The supplied research summary reports substantial complexes-of-groups and nonpositively curved combination theory in many cases, but not a fully general extension under a specified common set of hypotheses.
\subsection{Short English statement}
Find a higher-dimensional complexes-of-groups analogue of the Bestvina--Feighn combination theorem, with explicit conditions ensuring developability and hyperbolicity and recovering the graph case.
\subsection{Sources}
\begin{itemize}
\item Local UnsolvedMath registry, permanent ID \(6200090\), exact source code \texttt{catalogue entry}, ``Boundaries of Groups and Kleinian Groups --- Problem 90.'' Proposed year 2005; no proposer supplied.
\item Kapovich, \emph{Problems on Boundaries of Groups and Kleinian Groups} (2005), Problem 90, recorded PDF page 22; \url{https://www.math.ucdavis.edu/~kapovich/EPR/problems.pdf}. Recorded access 2026-07-29; pdf-text-normalized extraction.
\item The statement names the Bestvina--Feighn combination theorem without complete bibliographic details. The supplied PARTIAL-PROGRESS summary additionally names Haefliger, Corson, and Bridson for complexes-of-groups/combination theory, without precise references or theorem statements.
\end{itemize}
\end{document}
